%%%%%%%%%%%%%%%%%%%%%%%%%%%%%%%%%%%%%%%%%%%%%
\section{Mixing and linear response}
\label{Sec:Linear}
%%%%%%%%%%%%%%%%%%%%%%%%%%%%%%%%%%%%%%%%%%%%%

In this section we first review the mixing phenomenon in the absence of self-gravity and discuss the rate at which the observables introduced in the previous section decay. Next, we linearize the Vlasov-Poisson system about a self-consistent steady state and formulate, in terms of the action-angle variables of the latter, the criterion for the existence of oscillating modes. This leads to the gain $\lambda(\omega)$ of the self-consistent loop mentioned in the introduction, whose behavior at the edge of the frequency band is discussed at the end of the section. The derivations in this section are formal; for rigorous statements we refer the reader to Refs.~\cite{cM19,pRoS20} regarding the mixing property and to Refs.~\cite{mHgRcS22,Kunze-Book,mHgRmScS25,mMpRhV26} regarding the linearized system.


\subsection{Mixing in the absence of self-gravity}
\label{SubSec:FreeMixing}

When the self-gravity of the gas is neglected, $\Phi = \Phi_{ext}$ is independent of time, and the solution of Eq.~(\ref{Eq:VlasovAA}) is
\begin{equation}
\mathcal{F}(t,Q,J,L) = F(Q - \Omega(J,L)t,J,L),
\label{Eq:FreeSolution}
\end{equation}
where the function $F$, which is $2\pi$-periodic in its first argument, describes the initial DF. Expanding $F$ in a Fourier series with respect to the angle variable,
\begin{equation}
F(Q,J,L) = \sum\limits_{n\in\Integer} \hat{F}_n(J,L) e^{inQ},\qquad
\hat{F}_n(J,L) := \frac{1}{2\pi}\int\limits_0^{2\pi} F(Q,J,L) e^{-inQ} dQ,
\label{Eq:FourierF}
\end{equation}
and introducing Eq.~(\ref{Eq:FreeSolution}) into Eq.~(\ref{Eq:hk}) one obtains
\begin{equation}
h_n(t) = 16\pi^3\int\limits_0^\infty\int\limits_0^\infty \hat{F}_n(J,L) B(J,L) e^{-in\Omega(J,L)t}\, dJ\, L dL.
\label{Eq:hnFree}
\end{equation}
Hence, $h_0$ is constant in time, whereas for $n\neq 0$ the function $h_n(t)$ is a superposition of oscillations whose frequencies $n\Omega(J,L)$ vary from one orbit to the other. Provided that the gradient of $\Omega$ with respect to $(J,L)$ is different from zero for almost all $(J,L)$ (in the reduced model, provided that $\partial\Omega/\partial J\neq 0$ for almost all $J$), it follows from the Riemann-Lebesgue lemma and its generalizations (cf. Refs.~\cite{cM19,pRoS20}) that $h_n(t)\to 0$ for $t\to \infty$ and all $n\neq 0$. As a consequence, the observables~(\ref{Eq:NphiSum}) converge to the ones which belong to the angle-averaged DF $\hat{F}_0(J,L)$, which is a steady state according to the remark below Eq.~(\ref{Eq:VlasovAA}). This is the phase space mixing property. For the isochrone and the Kepler potentials, Eq.~(\ref{Eq:OmegaIsochrone}) gives $\partial\Omega/\partial J = -3[J + c(L)]^{-4} < 0$, such that the non-degeneracy condition is satisfied.

The rate at which $h_n(t)$ decays is not universal. It depends on the regularity of the function $\hat{F}_n B$ and on its behavior at the boundary of its support. If this function is smooth and vanishes together with all its derivatives at the boundary, repeated integration by parts in Eq.~(\ref{Eq:hnFree}) shows that $h_n(t)$ decays faster than any inverse power of $t$. A different situation, which is the relevant one for the steady states considered in this article, arises in the reduced model when the DF has its support in an interval $0\leq J\leq J_{max}$ and vanishes like a power at the outer edge, such that
\begin{equation}
\hat{F}_n(J) B(J) = (J_{max} - J)^k g(J)
\label{Eq:EdgeBehavior}
\end{equation}
with an exponent $k > 0$ and a smooth function $g$ satisfying $g(J_{max})\neq 0$. Suppose that $\Omega'(J) < 0$ for all $0 < J\leq J_{max}$, as is the case in all the examples of this article, and that $g$ vanishes sufficiently fast for $J\to 0$. Then, the behavior of the integral in Eq.~(\ref{Eq:hnFree}) for large times is determined by the endpoint $J = J_{max}$ (see, for instance, chapter~6 in Ref.~\cite{BenderOrszag-Book}), and for $n\geq 1$ one finds
\begin{equation}
h_n(t)\simeq 16\pi^3 L_0\,\Gamma(k+1)\, g(J_{max})\frac{e^{-in\Omega_{min} t}}{\left[ i n |\Omega'(J_{max})| t \right]^{k+1}},\qquad t\to \infty,
\label{Eq:Tail}
\end{equation}
where $\Omega_{min} := \Omega(J_{max})$ is the frequency of the outermost orbit and $\Gamma$ denotes the Gamma function. Therefore, an edge at which the DF vanishes with the exponent $k$ gives rise to an algebraic tail which is proportional to $t^{-(k+1)}$ and which oscillates with the frequency $n\Omega_{min}$: the less regular the DF is at the edge, the slower is the decay. Algebraic decay laws in related settings have been obtained in Refs.~\cite{jBaOyY11,mHgRmScS24}.

When the angular momenta of the gas particles are distributed, the integration over $L$ in Eq.~(\ref{Eq:hnFree}) provides an additional source of decay, since $\Omega$ depends on $L$ as well. For the isochrone and the Kepler potentials, Eq.~(\ref{Eq:OmegaIsochrone}) yields $\partial\Omega/\partial L = c'(L)\,\partial\Omega/\partial J$ with $1/2 < c'(L)\leq 1$, such that a dispersion in $L$ is almost as efficient for the mixing as a dispersion of the same magnitude in $J$. This effect is analyzed in section~\ref{Sec:Dispersion}.


\subsection{Linearization about a self-consistent steady state}
\label{SubSec:Linearization}

Let $(\mathcal{F}_{eq},\Phi_{eq})$ be a self-consistent steady state, and let $(Q,J)$ be the action-angle variables which belong to the potential $\Phi_{eq}$, such that $\mathcal{F}_{eq}$ is a function of $(J,L)$ only. We write
\begin{equation}
\mathcal{F} = \mathcal{F}_{eq} + \delta\mathcal{F},\qquad
\Phi = \Phi_{eq} + \delta\Phi,
\end{equation}
with $\delta\Phi$ the perturbation of the potential defined in Eq.~(\ref{Eq:deltaPhi}). The motion of the gas particles is determined by the Hamiltonian $H = E(J,L) + \delta\Phi(t,r(Q,J,L))$, where $r(Q,J,L)$ denotes the radius as a function of the action-angle variables. Since for each fixed $L$ the variables $(Q,J)$ are canonical, the Vlasov equation~(\ref{Eq:VlasovSph}) is equivalent to
\begin{equation}
\frac{\partial\delta\mathcal{F}}{\partial t} + \Omega\frac{\partial\delta\mathcal{F}}{\partial Q} - \frac{\partial\mathcal{F}_{eq}}{\partial J}\frac{\partial\delta\Phi}{\partial Q}
 = \frac{\partial\delta\Phi}{\partial Q}\frac{\partial\delta\mathcal{F}}{\partial J} - \frac{\partial\delta\Phi}{\partial J}\frac{\partial\delta\mathcal{F}}{\partial Q},
\label{Eq:VlasovPert}
\end{equation}
where no approximation has been made so far. The right-hand side is quadratic in the perturbation, and neglecting it yields the linearized equation
\begin{equation}
\frac{\partial\delta\mathcal{F}}{\partial t} + \Omega(J,L)\frac{\partial\delta\mathcal{F}}{\partial Q} - \frac{\partial\mathcal{F}_{eq}}{\partial J}(J,L)\frac{\partial\delta\Phi}{\partial Q} = 0.
\label{Eq:LinVlasov}
\end{equation}
The potential is determined by Eqs.~(\ref{Eq:PoissonSph},\ref{Eq:PhigasSol},\ref{Eq:PSVolumeForm}), which give
\begin{equation}
\delta\Phi(t,r) = -8\pi^2 G\int\limits_0^\infty\int\limits_0^\infty\int\limits_0^{2\pi} \frac{\delta\mathcal{F}(t,Q',J',L')}{\max\{ r,r(Q',J',L') \}}\, dQ'\, dJ'\, L' dL'.
\label{Eq:LinPoisson}
\end{equation}
The second term on the left-hand side of Eq.~(\ref{Eq:LinVlasov}) describes the transport along the orbits of the steady state discussed in the previous subsection. The third term describes the reaction of the steady state to the perturbation of the potential: the latter changes the action variables of the gas particles at the rate $\dot{J} = -\partial\delta\Phi/\partial Q$, and this displaces mass along the gradient of $\mathcal{F}_{eq}$. Through this term, which is absent when self-gravity is neglected, all the orbits are coupled to each other.

Since the support of $\mathcal{F}_{eq}$ is bounded, the frequencies $\Omega(J,L)$ of the orbits which are populated by the steady state fill an interval $[\Omega_{min},\Omega_{max}]$ with $0 < \Omega_{min} < \Omega_{max}$, which will be referred to as the frequency band in what follows. Without the third term in Eq.~(\ref{Eq:LinVlasov}), the $n$-th Fourier coefficient of $\delta\mathcal{F}$ with respect to $Q$ oscillates with the frequencies $n\Omega(J,L)$, which lie in the interval $n[\Omega_{min},\Omega_{max}]$. These intervals constitute the essential spectrum of the linearized system, which is not altered by the self-gravity of the perturbation~\cite{mHgRcS22,mHgRmScS25}.


\subsection{Oscillating modes and the gain of the self-consistent loop}
\label{SubSec:Gain}

We look for solutions of Eqs.~(\ref{Eq:LinVlasov},\ref{Eq:LinPoisson}) of the form
\begin{equation}
\delta\mathcal{F}(t,Q,J,L) = a(Q,J,L) e^{-i\omega t} + c.c.,\qquad
\delta\Phi(t,r) = \phi(r) e^{-i\omega t} + c.c.,
\label{Eq:ModeAnsatz}
\end{equation}
with a real frequency $\omega > 0$, where $c.c.$ denotes the complex conjugate of the preceding term. Such a solution describes a perturbation which oscillates without being damped, and it will be called an oscillating mode in what follows. We expand
\begin{equation}
a(Q,J,L) = \sum\limits_{n\in\Integer} a_n(J,L) e^{inQ},\qquad
\phi(r(Q,J,L)) = \sum\limits_{n\in\Integer} \phi_n(J,L) e^{inQ}.
\end{equation}
Since the radius is the same at the two points of an orbit with angle variables $Q$ and $2\pi - Q$, the coefficients of the potential satisfy $\phi_{-n} = \phi_n$ and are given by
\begin{equation}
\phi_n(J,L) = \frac{1}{2\pi}\int\limits_0^{2\pi} \phi(r(Q,J,L))\cos(nQ)\, dQ.
\label{Eq:phin}
\end{equation}
Equation~(\ref{Eq:LinVlasov}) yields
\begin{equation}
\left[ n\Omega(J,L) - \omega \right] a_n(J,L) = n\frac{\partial\mathcal{F}_{eq}}{\partial J}(J,L)\,\phi_n(J,L),\qquad n\in\Integer,
\label{Eq:ModeEq}
\end{equation}
which implies that $a_0 = 0$. From now on we assume that $\omega$ lies below the frequency band, that is $0 < \omega < \Omega_{min}$.\footnote{When $\Omega_{max} < 2\Omega_{min}$ there are further gaps between the intervals $n[\Omega_{min},\Omega_{max}]$. We do not consider frequencies lying in these gaps, since in all our examples the collective oscillation emerges below the band.} In this case none of the gas particles is in resonance with the perturbation, in the sense that $n\Omega(J,L)\neq \omega$ for all $n\in\Integer$ and all the orbits populated by the steady state, and Eq.~(\ref{Eq:ModeEq}) can be solved for $a_n$. By the same symmetry argument as above, the potential generated by the perturbation depends only on the combinations $b_n := a_n + a_{-n}$ with $n\geq 1$, for which Eq.~(\ref{Eq:ModeEq}) gives
\begin{equation}
b_n(J,L) = R_n(J,L;\omega)\,\phi_n(J,L),\qquad
R_n(J,L;\omega) := \frac{2n^2\Omega(J,L)}{n^2\Omega(J,L)^2 - \omega^2}\frac{\partial\mathcal{F}_{eq}}{\partial J}(J,L).
\label{Eq:Rn}
\end{equation}
On the other hand, introducing the ansatz~(\ref{Eq:ModeAnsatz}) into Eq.~(\ref{Eq:LinPoisson}) and using Eq.~(\ref{Eq:phin}) one obtains
\begin{equation}
\phi_n(J,L) = \sum\limits_{n'=1}^\infty\int\limits_0^\infty\int\limits_0^\infty M_{nn'}(J,L;J',L')\, b_{n'}(J',L')\, dJ'\, L' dL',
\label{Eq:PoissonHarm}
\end{equation}
with the kernel
\begin{equation}
M_{nn'}(J,L;J',L') := -4\pi G\int\limits_0^{2\pi}\int\limits_0^{2\pi} \frac{\cos(nQ)\cos(n'Q')}{\max\{ r(Q,J,L),r(Q',J',L') \}}\, dQ\, dQ'.
\label{Eq:Mnn}
\end{equation}
Equations~(\ref{Eq:Rn},\ref{Eq:PoissonHarm}) can be written in the compact form
\begin{equation}
\phi = \mathcal{K}(\omega)\phi,\qquad
\mathcal{K}(\omega) := M R(\omega),
\label{Eq:Loop}
\end{equation}
where $\phi = (\phi_n)_{n\geq 1}$, $R(\omega)$ denotes the multiplication with the functions $R_n(\cdot\,;\omega)$, and $M$ is the integral operator defined by the right-hand side of Eq.~(\ref{Eq:PoissonHarm}). The operator $\mathcal{K}(\omega)$ maps a perturbation of the potential with frequency $\omega$ to the potential which is generated by the response of the gas to this perturbation. Therefore, an oscillating mode with frequency $\omega$ exists if and only if one is an eigenvalue of $\mathcal{K}(\omega)$. In the reduced model, Eqs.~(\ref{Eq:phin}--\ref{Eq:Mnn}) hold with $L = L' = L_0$, the function $\mathcal{F}_{eq}$ replaced by $\mathcal{F}_{0,eq}(J)$, and the integral over $L'$ by the factor $L_0$.

The operators $M$ and $R(\omega)$ have the following properties. First, with respect to the inner product
\begin{equation}
\langle b,c \rangle := \sum\limits_{n=1}^\infty\int\limits_0^\infty\int\limits_0^\infty b_n(J,L) c_n(J,L)\, dJ\, L dL,
\label{Eq:InnerProduct}
\end{equation}
the operator $M$ is symmetric and negative, since for real-valued $b$ one finds
\begin{equation}
-\langle b,M b \rangle = \frac{1}{64\pi^4 G}\int\limits_{\Real^3} |\nabla\delta\Phi_b|^2 d^3 x\geq 0,
\label{Eq:MNegative}
\end{equation}
where $\delta\Phi_b$ denotes the potential which is obtained from Eq.~(\ref{Eq:LinPoisson}) when $\delta\mathcal{F}$ is replaced with $\sum_{n\geq 1} b_n(J,L)\cos(nQ)$. Second, suppose that $\partial\mathcal{F}_{eq}/\partial J < 0$ in the interior of the support of the steady state. Since $\partial\mathcal{F}_{eq}/\partial J = \Omega\,\partial\psi/\partial E$, this means that the function $\psi$ in Eq.~(\ref{Eq:SteadyState}) decreases with the energy, which is Antonov's stability condition~\cite{BinneyTremaine-Book,mHgRmScS25}, and it holds for all the steady states considered in this article. Then, $R_n(J,L;\omega) < 0$ for all $0\leq \omega < \Omega_{min}$, and $\mathcal{K}(\omega) = (-M)|R(\omega)|$ has the same nonvanishing eigenvalues as the symmetric and nonnegative operator
\begin{equation}
S(\omega) := |R(\omega)|^{1/2} (-M) |R(\omega)|^{1/2}.
\label{Eq:SDef}
\end{equation}
We denote by $\lambda(\omega)$ the largest eigenvalue of $S(\omega)$ and call it the gain of the self-consistent loop. Since $|R_n(J,L;\omega)|$ increases with $\omega$, the function $\lambda(\omega)$ is nondecreasing on the interval $[0,\Omega_{min})$, and we define its value at the edge of the band by
\begin{equation}
\lambda_{edge} := \lim\limits_{\omega\to \Omega_{min}} \lambda(\omega)\in (0,\infty].
\label{Eq:LambdaEdge}
\end{equation}
Third, the same construction applies to solutions which grow exponentially in time, $\delta\mathcal{F}\propto e^{\sigma t}$ with $\sigma > 0$, for which $\omega^2$ is replaced with $-\sigma^2$ in Eq.~(\ref{Eq:Rn}). In this case $|R_n|$ decreases to zero when $\sigma$ increases from zero to infinity, such that a growing solution exists if and only if $\lambda(0) > 1$. Hence, the linear stability of the steady state with respect to spherically symmetric perturbations requires that $\lambda(0)\leq 1$. For all the steady states of the reduced model which are analyzed in section~\ref{Sec:FixedL} we find $\lambda(0) < 0.6$.

Combining these properties, we arrive at the following criterion for a linearly stable steady state with $\lambda(0) < 1$. If $\lambda_{edge} > 1$, there exists an oscillating mode whose frequency $\omega_d$ lies below the frequency band and is determined by the equation
\begin{equation}
\lambda(\omega_d) = 1,\qquad 0 < \omega_d < \Omega_{min}.
\label{Eq:ModeFrequency}
\end{equation}
If $\lambda_{edge} < 1$, there are no oscillating modes with frequencies below the band. Further modes would appear if other eigenvalues of $S(\omega)$ exceeded one. This does not occur in our examples, in which the second largest eigenvalue remains smaller than $0.26$ up to the edge of the band. Note that $1/\lambda_{edge}$ is the factor by which the perturbation of the potential in Eq.~(\ref{Eq:LinVlasov}) would have to be multiplied for an oscillating mode to appear at the edge of the band; in this sense, $\lambda_{edge}$ measures the strength of the self-gravity of the perturbation relative to the one that is required to prevent its mixing.

This criterion is a formulation in terms of action-angle variables of the Birman-Schwinger principle of Refs.~\cite{sM90,mHgRcS22,Kunze-Book,mHgRmScS25}. In fact, $\mu > 0$ is an eigenvalue of $\mathcal{K}(\omega)$ if and only if the linearized system in which $\delta\Phi$ is multiplied by the factor $1/\mu$ possesses an oscillating mode with frequency $\omega$, and this is the property which characterizes the eigenvalues of the Birman-Schwinger operator of Ref.~\cite{mHgRmScS25} with spectral parameter $\omega^2$. Therefore, $\lambda(\omega)$ coincides with the norm of the operator of Mathur defined in that reference, and the interval $0 < \omega < \Omega_{min}$ corresponds to the principal gap of the spectrum. In Refs.~\cite{mHgRcS22,Kunze-Book,mHgRmScS25} the criterion is used to prove the existence or the absence of oscillating modes for certain classes of steady states. In section~\ref{Sec:Numerics} we describe a numerical method for the computation of $\lambda(\omega)$, which allows us to evaluate the criterion for steady states of arbitrary mass. \EJ{As far as we could determine, this operator has not been evaluated numerically in the literature; Kunze's lecture notes~\cite{mK22} and Straub's thesis list it as an open task. Mathur's article~\cite{sM90} and Kunze's book~\cite{Kunze-Book} have to be checked before claiming this.}


\subsection{Behavior at the edge of the frequency band}
\label{SubSec:Edge}

The value of $\lambda_{edge}$ depends strongly on the way in which the DF of the steady state vanishes at the boundary of its support. To explain this, consider the reduced model and a steady state whose DF has a cutoff at the energy $E_0$ and behaves like
\begin{equation}
\mathcal{F}_{0,eq}\simeq \alpha (E_0 - E)^k,\qquad E\to E_0,
\label{Eq:PolytropeEdge}
\end{equation}
with constants $\alpha > 0$ and $k > 1/2$, as is the case for the polytropes considered in Ref.~\cite{mHgRmScS25}. Let $J_{max}$ be the action variable of the orbit with energy $E_0$, such that $\Omega_{min} = \Omega(J_{max})$ when $\Omega$ is a decreasing function of $J$. Then $\partial\mathcal{F}_{0,eq}/\partial J$ is proportional to $(J_{max} - J)^{k-1}$ near the edge, and Eq.~(\ref{Eq:Rn}) with $n = 1$ and $\omega = \Omega_{min}$ yields
\begin{equation}
|R_1(J;\Omega_{min})|\simeq \frac{1}{\Omega(J) - \Omega_{min}}\left| \frac{\partial\mathcal{F}_{0,eq}}{\partial J}(J) \right| \propto (J_{max} - J)^{k-2},\qquad J\to J_{max}.
\label{Eq:R1Edge}
\end{equation}
This function is integrable at $J = J_{max}$ if and only if $k > 1$, which leads to the following two cases.

If $k > 1$, $\lambda_{edge}$ is finite. Since $R(\omega)$ is proportional to the DF of the steady state, $\lambda_{edge}$ is proportional to $M_{gas}/M$ as long as this ratio is small enough for the contribution of $\Phi_{gas}$ to $\Phi_{eq}$ to be negligible. Therefore, for a gas of small mass one has $\lambda_{edge} < 1$ and there is no oscillating mode, in accordance with the damping result in Ref.~\cite{mHgRmScS25}. When the mass of the gas increases, $\lambda_{edge}$ grows, and an oscillating mode appears when it reaches the value one. The corresponding threshold mass is computed in section~\ref{Sec:FixedL}. Close to the edge one finds $\lambda_{edge} - \lambda(\omega)\propto (\Omega_{min} - \omega)^{\min\{ k-1,1 \}}$, with a logarithmic correction when $k = 2$. Hence, slightly above the threshold the frequency of the mode satisfies $\Omega_{min} - \omega_d\propto (\lambda_{edge} - 1)^{1/(k-1)}$ when $1 < k < 2$, such that the closer $k$ lies to one, the closer the frequency of the mode stays to the edge of the band.

If $1/2 < k\leq 1$, $\lambda(\omega)$ diverges for $\omega\to \Omega_{min}$, like $(\Omega_{min} - \omega)^{k-1}$ when $k < 1$ and like $|\log(\Omega_{min} - \omega)|$ when $k = 1$. Hence, $\lambda_{edge} = \infty$ and an oscillating mode exists for any mass of the gas, as proven in Ref.~\cite{mHgRmScS25} for small masses. However, since the factor which multiplies the divergent term is proportional to $M_{gas}/M$, Eq.~(\ref{Eq:ModeFrequency}) implies that, to leading order in the limit $M_{gas}/M\to 0$,
\begin{equation}
\Omega_{min} - \omega_d\propto \left( \frac{M_{gas}}{M} \right)^{\frac{1}{1-k}}\quad (k < 1),\qquad
\Omega_{min} - \omega_d\propto \exp\left( -C\frac{M}{M_{gas}} \right)\quad (k = 1),
\label{Eq:Binding}
\end{equation}
with a positive constant $C$. Therefore, when the mass of the gas is small the frequency of the mode lies extremely close to the edge of the band. Since a time of the order of $(\Omega_{min} - \omega_d)^{-1}$ is required to distinguish the mode from the response of the orbits close to the edge, which oscillates with a frequency close to $\Omega_{min}$ and decays algebraically (cf. Eq.~(\ref{Eq:Tail})), in this regime the mode cannot be observed in a time evolution of practical duration.

When the angular momenta of the gas particles are distributed, the support of the steady state in the $(J,L)$-plane is two-dimensional, and in general the minimum $\Omega_{min}$ of the frequency is attained only at isolated points of its boundary. The additional integration over $L$ in Eq.~(\ref{Eq:PoissonHarm}) then weakens the singularity of $R_1$ at the edge of the band, such that $\lambda_{edge}$ can be finite also for $k\leq 1$, in which case the oscillating modes of the reduced model disappear when the mass of the gas is small enough. This is analyzed in section~\ref{Sec:Dispersion}. \EJ{The precise statement depends on the family of steady states: for a cutoff in the energy and the isochrone or Kepler potential, $\Omega$ depends on $L$ at fixed energy only through the self-gravity. To be settled together with section~\ref{Sec:Dispersion}.}

Finally, we comment on the case $\lambda_{edge} < 1$, in which the linearized system has no oscillating mode below the band. As we show in section~\ref{Sec:FixedL}, in this case the perturbation decays. When the self-gravity of the gas is not too weak, the response to an initial perturbation is dominated at intermediate times by a damped oscillation proportional to $e^{-i\omega t - \gamma t}$ with a frequency $\omega$ inside the band and a damping rate $\gamma > 0$, which is followed at late times by an algebraic tail of the form~(\ref{Eq:Tail}). The damped oscillation is the analogue of Landau damping in a plasma~\cite{lL46}. It does not correspond to an eigenvalue of the linearized system; instead, $\omega - i\gamma$ is a zero of the analytic continuation of the function $\omega\mapsto \det[1 - \mathcal{K}(\omega)]$ from the upper to the lower half of the complex plane, see for instance Refs.~\cite{BinneyTremaine-Book,mW94,jFsP22}. We do not compute this continuation in this article. Instead, we determine $\omega$ and $\gamma$ by fitting the time series which are obtained from the numerical solution of the linearized equations and from the PIC code, as described in section~\ref{Sec:Numerics}.
